A language model response can open with harmless background and only turn into actionable harm hundreds of tokens later. Streaming monitors and activation verbalizers are almost always evaluated on whether text is \emph{already} unsafe. We ask whether a still-safe prefix carries enough signal to forecast that a clear violation will first appear within the next $H\in\{64,128,256,512\}$ tokens. We build a family-split benchmark of 3,520 safe prefixes from 880 prompt families of fresh \qwen generations, with LLM-annotated token-level onsets and censoring-aware labels. On it we compare text classifiers, activation probes, model self-report, an off-the-shelf Activation Oracle, and two readers trained on future-event labels: a direct question-answering reader and a reader with an explicit forecast representation (horizon queries and a timing head). Delayed onsets are partly forecastable. A linear probe on a 16-token activation history reaches \auroc 0.79 at $H{=}64$ and 0.90 at $H{=}512$, and it flags 56\% of unsafe trajectories at least 128 tokens before onset at a 5\% benign false-alarm rate. The off-the-shelf Activation Oracle answers ``No'' to about 99\% of prefixes and is near chance for imminent onsets (\auroc 0.54). The forecast representation does not beat the matched direct reader. It is tied at $H\le128$ and slightly worse at $H\ge256$ ($\Delta$\auroc $-0.03$ and $-0.02$, family-paired bootstrap CIs excluding zero). In this setting, future-event supervision rather than forecast architecture carries the gains, and a cheap activation-history probe is the baseline to beat.
